\subsection{Valuations on a field}

PROPOSITION 2. Let K be a (not necessarily commutative) field and v a valuation on K with values in Γ. Then:

\begin{enumerate}
\def\labelenumi{(\roman{enumi})}
\item
  For \(x \neq 0\) , \(v(x) \neq +\infty\) .
\item
  The set \(\mathbf{A}\) of \(x\in \mathbf{K}\) such that \(v(x)\geqslant 0\) is a subring of \(\mathbf{K}\) .
\item
  For all \(\alpha \geqslant 0\) in \(\Gamma\) , the set \(\mathrm{V}_{\alpha}\) (resp. \(\mathrm{V}_{\alpha}^{\prime}\) ) of \(x \in \mathbf{A}\) such that \(v(x) > \alpha\) (resp. \(v(x) \geqslant \alpha\) ) is a two-sided ideal of \(\mathbf{A}\) and every (left or right) ideal \(\neq (0)\) of \(\mathbf{A}\) contains one of the \(\mathrm{V}_{\alpha}^{\prime}\) .
\item
  The set \(\mathfrak{m}(\mathbf{A})\) of \(x \in \mathbf{A}\) such that \(v(x) > 0\) is the greatest ideal \(\neq \mathbf{A}\) in \(\mathbf{A}\) ; \(\mathbf{U}(\mathbf{A}) = \mathbf{A} - \mathfrak{m}(\mathbf{A})\) is the set of invertible elements of \(\mathbf{A}\) and \(\kappa(\mathbf{A}) = \mathbf{A} / \mathfrak{m}(\mathbf{A})\) is a (not necessarily commutative) field.
\item
  For all \(x \in \mathbf{K} - \mathbf{A}\) , \(x^{-1} \in \mathfrak{m}(\mathbf{A})\) .
\end{enumerate}

Assertion (i) follows from the fact that \(\overline{v}^{-1}(+\infty)\) is an ideal of K not equal to K. The verification of the fact that A is a ring and the \(V_{\alpha}\) and \(V_{\alpha}^{\prime}\) two-sided ideals is trivial by virtue of axioms ( \(VL_{I}\) ), ( \(VL_{II}\) ) and ( \(VL_{III}\) ). If a is a (left, for example) ideal of A and \(x \neq 0\) belongs to A, every \(y \in A\) such that \(v(y) \geqslant v(x)\) can be written y = zx where \(z = yx^{-1}\) , hence \(v(z) = v(y) - v(x) \geqslant 0\) and therefore \(z \in A\) ; in other words the left ideal Ax contains the \(V_{\alpha}^{\prime}\) for \(\alpha \geqslant v(x)\) . The set \(U(A) = A - m(A)\) is the set of \(x \in K\) such that \(v(x) = 0\) ; if \(x \in U(A)\) , then

\[
v (x ^ {- 1}) = - v (x) = 0,
\]

whence \(x^{-1} \in \mathrm{U}(\mathrm{A})\) ; conversely, if \(y \in \mathrm{A}\) is invertible in \(\mathrm{A}\) , then \(v(y) \geqslant 0\) , \(v(y^{-1}) \geqslant 0\) and \(v(y) + v(y^{-1}) = 0\) , whence \(v(y) = 0\) and \(y \in \mathrm{U}(\mathrm{A})\) ; this proves (iv) and (v) follows immediately from the definitions.

A (resp. m(A), κ(A)) is called the ring (resp. ideal, residue field) of the valuation v on K.

Clearly U(A) is the kernel of the homomorphism \(v: K^{*} \to \Gamma\) and the image \(v(\mathbf{K}^{*})\) under v of the multiplicative group \(K^{*}\) is a subgroup of the additive group \(\Gamma\) , called the order group or value group of v, which is therefore isomorphic to \(K^{*}/U(A)\) ; for \(x \in K\) , the element \(v(x)\) of \(\Gamma_{\infty}\) is sometimes called the valuation or order of x for v. Two valuations \(v, v'\) on K are called equivalent if they have the same ring.

PROPOSITION 3. For two valuations \(v, v'\) over \(a\) (not necessarily commutative) field \(K\) to be equivalent, it is necessary and sufficient that there exist an isomorphism \(\lambda\) of the ordered group \(v(K^*)\) onto the ordered group \(v'(K^*)\) such that \(v' = \lambda \circ v\) .

Suppose \(v\) and \(v'\) are equivalent; by hypothesis, the ring A of the valuation \(v\) being the same as that of \(v'\) , \(v\) and \(v'\) (restricted to K*) factor into homomorphisms \(K^* \to K^*/U(A) \xrightarrow{\mu} v(K^*)\) , \(K^* \to K^*/U(A) \xrightarrow{\nu} v'(K^*)\) , where \(\mu\) and \(\nu\) are isomorphisms; moreover, the set of positive elements of \(v(K^*)\) (resp. \(v'(K^*)\) ) is the image under \(\mu\) (resp. \(\nu\) ) of the set of classes mod. U(A) of elements \(\neq 0\) of m(A); we conclude that \(\lambda = \nu \circ \overline{\mu}^{-1}\) solves the problem, the converse being obvious.

Suppose now that K is a commutative field; then, for all valuations v on K, the ring A of the valuation v is a valuation ring of K in the sense of § 1, no. 2, Definition 2 (which justifies the terminology); this follows immediately from Proposition 2 (c) and § 1 no. 2, Theorem 1 (c). Conversely, recall that for every integral domain B whose field of fractions is K the relation of divisibility x\textbar y (equivalent to \(y \in Bx\) ) makes \(K^{*}\) into a preordered group, whose associated ordered group \(\Gamma_{B}\) is the quotient \(K^{*}/U(B)\) of \(K^{*}\) by the group \(U(B)\) of invertible elements of B, the positive elements of this group being those of \(B^{*}/U(B)\) (where \(B^{*} = B - \{0\}\) ); the mapping \(x \mapsto Bx\) defines, by passing to the quotient, an isomorphism of the ordered group \(K^{*}/U(B)\) onto the group (ordered by the relation \(\supset\) ) of non-zero principal fractional ideals of K (Algebra, Chapter VI, § 1, no. 5). The rings A with field of fractions K and for which the group \(\Gamma_{A} = K^{*}/U(A)\) is totally ordered are precisely the valuation rings of K (§ 1, no. 2, Theorem 1 (d)). If \(v_{A}\) denotes the canonical homomorphism of \(K^{*}\) onto \(\Gamma_{A}\) , it is immediate that \(v_{A}\) (extended by \(v_{A}(0) = +\infty\) ) is a valuation (called canonical) on K whose ring is A; every valuation equivalent to \(v_{A}\) may be written \(v = \sigma \circ v_{A}\) , where \(\sigma\) is an isomorphism of \(\Gamma_{A}\) onto a subgroup of the group where v takes its values (Proposition 3); \(\sigma \circ v_{A}\) is called the canonical factorization of v.

PROPOSITION 4. Let C be an integral domain, K its field of fractions, \(C^{*} = C - \{0\}\) and \(v: C \to \Gamma_{\infty}\) a valuation on C. Then there exists a unique valuation w on K which extends v and \(w(K^{*})\) is the subgroup of \(\Gamma\) generated by \(v(C^{*})\) .

By Theorem 2 of Algebra, Chapter I, § 2, no. 7, there exists a unique homomorphism \(w\) from \(\mathbf{K}^*\) to \(\Gamma\) which extends \(v \mid \mathbf{C}^*\) and \(w(\mathbf{K}^*)\) is generated by \(v(\mathbf{C}^*)\) . It remains to prove that \(w\) satisfies axiom \((\mathrm{VL}_{\mathrm{II}})\) . Then let \(x' \in \mathbf{K}^*\) , \(y \in \mathbf{K}^*\) be such that \(x + y \in \mathbf{K}^*\) ; there exists \(a \in \mathbf{C}^*\) such that \(ax \in \mathbf{C}^*\) and \(ay \in \mathbf{C}^*\) , whence \(a(x + y) \in \mathbf{C}^*\) . Since the restriction of \(w\) to \(\mathbf{C}^*\) satisfies \((\mathrm{VL}_{\mathrm{II}})\) ,

\[
w (a (x + y)) \geqslant \inf (w (a x), w (a y)).
\]

Eliminating \(w(a)\) from both sides, we obtain

\[
w (x + y) \geqslant \inf (w (x), w (y)).
\]

